\documentclass[11pt,a4paper,oneside]{book}

\usepackage[T1]{fontenc}
\usepackage[utf8]{inputenc}
\usepackage[english]{babel}
\usepackage{lmodern}
\usepackage[a4paper,inner=30mm,outer=27mm,top=28mm,bottom=30mm]{geometry}
\usepackage{amsmath,amssymb,amsthm,mathrsfs,mathtools,stmaryrd}
\usepackage{tikz,tikz-cd}
\usetikzlibrary{
  arrows,
  arrows.meta,
  calc,
  decorations.pathmorphing,
  matrix,
  positioning,
  shapes,
  shapes.geometric
}
\usepackage{graphicx,float,adjustbox,enumitem,multicol}
\usepackage{etoolbox}
\usepackage{xcolor}
\usepackage[hidelinks,hypertexnames=false,pdfstartview={},linktoc=page,
  pdftitle={SGA 3 - Complete English Translation and Native-Diagram Reconstruction},
  pdfauthor={SGA 3 English reconstruction project},
  pdfsubject={Cumulative English reader for SGA 3: Introduction, Exposes I-XXVI, Tome-I subject index, Tome-III guide, and terminal index},
  pdfkeywords={SGA 3, group schemes, English translation, native TeX diagrams}]{hyperref}

\setlength{\emergencystretch}{3em}
\raggedbottom

\theoremstyle{plain}
\newtheorem{theorem}{Theorem}[section]
\newtheorem{proposition}[theorem]{Proposition}
\newtheorem{corollary}[theorem]{Corollary}
\newtheorem{lemma}[theorem]{Lemma}
\newtheorem{scholium}[theorem]{Scholium}
\newtheorem{application}[theorem]{Application}
\newtheorem{complements}[theorem]{Complements}

\theoremstyle{definition}
\newtheorem{definition}[theorem]{Definition}
\newtheorem{example}[theorem]{Example}

\theoremstyle{remark}
\newtheorem{remark}[theorem]{Remark}

\providecommand{\Ga}{\mathbb G_a}
\providecommand{\Gm}{\mathbb G_m}
\providecommand{\Norm}{\operatorname{Norm}}
\providecommand{\Centr}{\operatorname{Centr}}
\providecommand{\Ker}{\operatorname{Ker}}
\providecommand{\Dom}{\operatorname{Dom}}
\providecommand{\rank}{\operatorname{rank}}
\providecommand{\Spec}{\operatorname{Spec}}
\providecommand{\Lie}{\operatorname{Lie}}
\providecommand{\intaut}{\operatorname{int}}
\providecommand{\Fib}{\operatorname{Fib}}
\providecommand{\Aut}{\operatorname{Aut}}
\providecommand{\Autext}{\operatorname{Autext}}
\providecommand{\rgred}{\operatorname{rgred}}
\providecommand{\rgss}{\operatorname{rgss}}
\providecommand{\GL}{\operatorname{GL}}
\providecommand{\Qp}{\mathbb{Q}_{+}}
\providecommand{\cR}{\mathcal{R}}
\providecommand{\VR}{\mathcal{V}}
\providecommand{\Nil}{\operatorname{Nil}}

\providecommand{\SGAParagraph}[1]{%
  \par\medskip\noindent\textbf{#1.}---\quad
}
\providecommand{\NDE}[1]{%
  \footnote{\textit{Editors' note.} #1}%
}
\newenvironment{sgastatement}[1]
  {\par\medskip\noindent\textbf{#1.}---\ \itshape}
  {\par\medskip}
\providecommand{\SGAQED}{\hfill\(\square\)\par\smallskip}
\providecommand{\Idx}[2]{%
  \par\noindent\hangindent=1.2em\hangafter=1
  #1\if\relax\detokenize{#2}\relax\else,\ #2\fi\par
}
\providecommand{\SGASection}[2]{%
  \section*{#1.\ #2}%
  \addcontentsline{toc}{section}{#1.\ #2}%
}
\providecommand{\SGASubsection}[2]{%
  \subsection*{#1.\ #2}%
  \addcontentsline{toc}{subsection}{#1.\ #2}%
}
\providecommand{\SGAStatement}[2]{%
  \par\medskip\noindent\textbf{#1.}---\ {\itshape #2\par}\medskip
}
\providecommand{\SGAProof}[1]{%
  \par\smallskip\noindent\textit{#1.}\quad
}

\makeatletter
\providecommand{\SGARefTarget}[1]{\phantomsection\label{#1}}
\providecommand{\SGAInlineRefTarget}[1]{%
  \begingroup
    \def\@currentHref{#1}%
    \label{#1}%
  \endgroup}
\providecommand{\SGAInlineEquationTarget}[1]{%
  \begingroup
    \def\@currentHref{#1}%
    \ltx@label{#1}%
  \endgroup}
\let\SGA@PendingSectionTarget\@empty
\providecommand{\SGANextSectionTarget}[1]{%
  \gdef\SGA@PendingSectionTarget{#1}}
\providecommand{\SGA@UsePendingSectionTarget}{%
  \ifx\SGA@PendingSectionTarget\@empty\else
    \expandafter\SGAInlineRefTarget\expandafter{%
      \SGA@PendingSectionTarget}%
    \global\let\SGA@PendingSectionTarget\@empty
  \fi}
\renewcommand{\SGASection}[2]{%
  \section*{\SGA@UsePendingSectionTarget#1.\ #2}%
  \addcontentsline{toc}{section}{#1.\ #2}%
}
\renewcommand{\SGASubsection}[2]{%
  \subsection*{\SGA@UsePendingSectionTarget#1.\ #2}%
  \addcontentsline{toc}{subsection}{#1.\ #2}%
}
\pdfstringdefDisableCommands{%
  \def\SGAInlineRefTarget#1{}%
  \def\SGAInlineEquationTarget#1{}%
  \def\SGANextSectionTarget#1{}%
}
\providecommand{\SGARefLink}[2]{%
  \@ifundefined{r@#1}{#2}{\hyperref[#1]{#2}}%
}
\providecommand{\GuideChapter}[2]{%
  \@dottedtocline{0}{0em}{0em}{#1}{#2}}
\providecommand{\GuideChapterSplit}[3]{%
  \par\noindent\textbf{#1}\par\nobreak
  \@dottedtocline{0}{2em}{0em}{#2}{#3}}
\providecommand{\GuideSection}[2]{%
  \@dottedtocline{1}{1.45em}{2.25em}{#1}{#2}}
\newcommand{\SGAInputRoot}[1]{%
  \def\input@path{{inputs/#1/}}%
}
\newcommand{\SGATagsLeft}{%
  \tagsleft@true
  \let\veqno\@@leqno
}
\newcommand{\SGATagsRight}{%
  \tagsleft@false
  \let\veqno\@@eqno
}
\makeatother

\newcommand{\SGAStartExpose}[1]{%
  \clearpage
  \setcounter{section}{0}%
  \setcounter{subsection}{0}%
  \setcounter{theorem}{0}%
  \setcounter{footnote}{0}%
  \setcounter{figure}{0}%
  \SGAInputRoot{#1}%
}

\title{Seminar on Algebraic Geometry 3\\
  \large Group Schemes\\[0.75em]
  \normalsize Complete English Translation and Native-Diagram
  Reconstruction\\
  \small Introduction; Expos\'es I--XXVI; Tome-I subject index;\\
  Tome-III mathematical guide and terminal index}
\author{SGA 3 English reconstruction project}
\date{29 July 2026}

\begin{document}
\begingroup\small
Rights in the original works and in any protected contributions to these derivative translations and transcriptions remain with their respective rights holders. Anything in these editions that the publisher holds rights to is dedicated to the public domain under CC0 1.0 (https://creativecommons.org/publicdomain/zero/1.0/). Existing third-party notices and any permissions already granted remain applicable.

These editions are made available free of charge for non-commercial, purely educational purposes, on the presumption that the respective rights holders would not object to these translations and transcriptions being freely available alongside the French originals on \href{https://www.numdam.org/journals/PMIHES/}{NUMDAM}. This presumption is not a statement that permission has been obtained; free access to the originals does not itself grant permission to translate, adapt or republish them.

If you hold rights in any material presented here and object to its availability, please \href{https://github.com/KokunoYumeto/sga-en/issues}{open a GitHub issue}. The publisher will take down the affected material at the rights holder's request.

The additional English translation is AI-translated. Corrections, editing and typesetting are AI-produced. The publisher did not translate the text. Human source authorship remains credited.
\par\endgroup\clearpage


\frontmatter
\maketitle

\chapter*{Sources and authorship}
The original SGA volumes retain their authors, seminar contributors and editors, named in the corresponding texts. The collection includes work of Alexander Grothendieck, Michèle Raynaud, Michel Raynaud, Pierre Deligne, Nicholas Katz, Michel Demazure, Michael Artin, Jean-Louis Verdier, Luc Illusie and their collaborators; this list does not replace the detailed credits in each volume.

SGA 1 retains the French retypesetting associated with Bas Edixhoven’s volunteer project, the Société mathématique de France and Michel Raynaud’s notes, from \href{https://arxiv.org/abs/math/0206203v2}{arXiv math/0206203v2}.

The SGA 2 and SGA 3 English sources acknowledge Jacob C. Reinhold’s \href{https://github.com/jcreinhold/sga/tree/e7a259f3f8608ad3edf9bf6eead3fd504dd2d23e}{English translation}, snapshot e7a259f3f8608ad3edf9bf6eead3fd504dd2d23e, used for comparison and drafting. His stated CC BY 4.0 notice applies to that contribution and is preserved; it does not license the underlying French works or every contribution here. Inherited English wording remains attributed to its human source. The SGA 3 French reference is the Philippe Gille–Patrick Polo re-edition.

SGA 4½ preserves Daniel Miller’s \href{https://github.com/dkmiller/sga4.5}{LaTeX typesetting} and its original typesetter’s note. It is not attributed to AI.

Additional English translation is AI-translated; subsequent corrections, editing, reconstruction and typesetting are AI-produced. The publisher did not translate the text. The source-specific human credits and genuine third-party notices remain applicable; attribution does not establish permission or endorsement.


\tableofcontents

\mainmatter

\SGAInputRoot{INTRO}
\input{body}

\SGAStartExpose{I}
\input{body}

\makeatletter
\providecommand{\SGAThreeNativeInput}[2][0.92\textwidth]{%
  \begin{adjustbox}{max width=#1,center,varwidth=\maxdimen}%
    \input{#2}%
  \end{adjustbox}%
}
\providecommand{\SGAThreeNativeDiagram}[5][0.92\textwidth]{%
  \begin{figure}[H]
    \centering
    \SGAThreeNativeInput[#1]{#2}%
    \caption{#3}
    \label{#4}
    \par\smallskip
    {\footnotesize\itshape Native \TeX{} reconstruction; source locator: #5.}
  \end{figure}%
}
\makeatother

\SGAStartExpose{II}
\input{body}

\SGAStartExpose{III}
\input{body}

\SGAStartExpose{IV}
\input{body}

\SGAStartExpose{V}
\makeatletter
\renewcommand{\SGAThreeNativeInput}[1]{%
  \filename@parse{#1}%
  \edef\SGAThreeNativeFile{%
    native_diagrams/exp5/\filename@base.\filename@ext.tex%
  }%
  \input{\SGAThreeNativeFile}%
}
\makeatother
\providecommand{\SGAThreeDiagram}[5][0.92\textwidth]{}
\renewcommand{\SGAThreeDiagram}[5][0.92\textwidth]{%
  \begin{figure}[H]
    \centering
    \resizebox{#1}{!}{\SGAThreeNativeInput{#2}}%
    \caption{#3}
    \label{#4}
    \par\smallskip
    {\footnotesize\itshape Native \LaTeX{} reconstruction from the
      corrected Polo--Gille diagram; source locator: #5.}
  \end{figure}%
}
\input{body}

\SGAStartExpose{VI}
\input{body}

\SGAStartExpose{VII}
\input{body}

\clearpage
\chapter*{Index --- Tome I}
\addcontentsline{toc}{chapter}{Index --- Tome I}
\markboth{Index --- Tome I}{Index --- Tome I}
\noindent\emph{Page numbers refer to the Polo--Gille re-edition.}
\medskip
\begin{multicols}{2}
\small
\setlength{\parskip}{1.2pt}
\interlinepenalty=10000
\SGAInputRoot{TOMEI_INDEX}
\input{body}
\end{multicols}

\newgeometry{margin=30mm}
\SGATagsLeft

\SGAStartExpose{VIII}
\input{body}

\SGAStartExpose{IX}
\input{body}

\SGAStartExpose{X}
\input{body}

\SGAStartExpose{XI}
\input{body}

\SGAStartExpose{XII}
\chapter*{Expos\'e XII\\
  Maximal Tori, the Weyl Group, Cartan Subgroups, and the Reductive Center}
\addcontentsline{toc}{chapter}{Expos\'e XII. Maximal Tori, the Weyl Group,
  Cartan Subgroups, and the Reductive Center}
\begin{center}
  \textit{by A. Grothendieck}
\end{center}
\input{body}

\SGAStartExpose{XIII}
\input{body}

\SGAStartExpose{XIV}
\chapter*{Expos\'e XIV\\Regular Elements: Continuation\\
Applications to Algebraic Groups}
\addcontentsline{toc}{chapter}{Expos\'e XIV. Regular Elements: Continuation}
\begin{center}
  \textit{by A. Grothendieck, with an appendix by J.-P. Serre}
\end{center}
\input{body}

\SGAStartExpose{XV}
\input{body}

\SGAStartExpose{XVI}
\input{body}

\SGAStartExpose{XVII}
\input{body}

\SGAStartExpose{XVIII}
\input{body}

\clearpage
\SGAInputRoot{GUIDE}
\input{body}

\SGAStartExpose{XIX}
\input{body}

\SGAStartExpose{XX}
\chapter*{Expos\'e XX\\Reductive Groups of Semisimple Rank 1}
\addcontentsline{toc}{chapter}{Expos\'e XX. Reductive Groups of Semisimple Rank 1}
\begin{center}
  \textit{by M. Demazure}
\end{center}
\input{body}

\SGAStartExpose{XXI}
\counterwithin{theorem}{subsection}
\renewcommand{\thesubsection}{\thesection.\arabic{subsection}}
\input{body}
\counterwithin{theorem}{section}

\SGAStartExpose{XXII}
\input{body}

\SGAStartExpose{XXIII}
\renewcommand{\SGAProof}{%
  \par\smallskip\noindent\textit{Proof.}\quad
}
\input{body}

\SGAStartExpose{XXIV}
\input{body}

\SGAStartExpose{XXV}
\input{body}

\SGAStartExpose{XXVI}
\providecommand{\rad}{\operatorname{rad}}
\providecommand{\Transt}{\underline{\operatorname{Transt}}}
\providecommand{\Tor}{\underline{\operatorname{Tor}}}
\providecommand{\Lev}{\underline{\operatorname{Lev}}}
\providecommand{\Par}{\underline{\operatorname{Par}}}
\providecommand{\Crit}{\underline{\operatorname{Crit}}}
\providecommand{\Bor}{\underline{\operatorname{Bor}}}
\providecommand{\Gen}{\underline{\operatorname{Gen}}}
\providecommand{\Opp}{\underline{\operatorname{Opp}}}
\providecommand{\Stand}{\underline{\operatorname{Stand}}}
\providecommand{\TypeStand}{\underline{\operatorname{TypeStand}}}
\providecommand{\Of}{\operatorname{Of}}
\providecommand{\Int}{\operatorname{int}}
\renewcommand{\Centr}{\underline{\operatorname{Centr}}}
\renewcommand{\Norm}{\underline{\operatorname{Norm}}}
\input{body}

\clearpage
\SGAInputRoot{TERMINAL_INDEX}
\input{body}

\clearpage
\chapter*{Cumulative boundary}
\addcontentsline{toc}{chapter}{Cumulative boundary}
This reader ends with the terminal Tome-III index.  It contains no
material beyond SGA~3 and makes no archive, publication, public-license,
or public-readback claim.

\end{document}
